\documentclass[11pt,letterpaper]{article}

\usepackage[margin=1in]{geometry}
\usepackage{graphicx}
\usepackage{booktabs}
\usepackage{amsmath}
\usepackage{hyperref}
\usepackage{xcolor}
\usepackage{caption}
\usepackage{float}

\hypersetup{colorlinks=true, linkcolor=blue, citecolor=blue, urlcolor=blue}

\newcommand{\pp}{\mbox{$p$+$p$}}
\newcommand{\sqs}{\mbox{$\sqrt{s}$}}
\newcommand{\pT}{\ensuremath{p_{\text{T}}}}
\newcommand{\zvtx}{\ensuremath{z_{\mathrm{vtx}}}}
\newcommand{\Lint}{\ensuremath{\mathcal{L}_{\text{int}}}}
\newcommand{\veff}{\ensuremath{\varepsilon_{\text{vtx+MBD}}}}

\graphicspath{{../../plotting/figures/lumi_convention/}}

\title{%
  \vspace{-1cm}
  {\large \textbf{sPHENIX Internal Note}}\\[0.5cm]
  {\LARGE \textbf{Luminosity Convention for PPG12:\\
  L(60cm) $\rightarrow$ L(noVtx) with Vertex Efficiency Fix}}\\[0.3cm]
  {\large Isolated Photon Cross-Section in \pp{} at $\sqs = 200$~GeV}
}
\author{sPHENIX Collaboration}
\date{\today}

\begin{document}

\maketitle

\begin{abstract}
We switch the PPG12 cross-section normalization from the
60~cm-vertex-cut luminosity $\Lint^{60\text{cm}}$ (with a vertex
efficiency whose denominator is truth photons within
$|\zvtx^{\text{truth}}|<60$~cm) to the no-vertex-cut luminosity
$\Lint^{\text{noVtx}}$ (with an efficiency denominator that includes
ALL truth photons). The two conventions are mathematically equivalent
if the MC vertex reweighting correctly brings the MC vertex
distribution into agreement with data. In the old production
(June 2025), they disagreed by a factor of $\sim 3$ because the MC
vertex reweighting used a deprecated static fallback file with a
\texttt{vertex\_weight = 1.0} fallback that kept raw MC events at
large $|\zvtx|$. We diagnosed the bug, removed the fallback entirely,
and hardened the on-the-fly reweight with rebinning, smoothing,
linear interpolation, and a \texttt{vertex\_weight = 0.0} fallback.
A standalone simulation using the existing first-pass vtxscan files
shows that after the fix the two conventions agree to $+0.32\%$
flat across all \pT{} bins, well below any relevant systematic.
\end{abstract}

\tableofcontents
\newpage

%% ============================================================
\section{Motivation}
\label{sec:motivation}

The PPG12 cross-section factorizes as:
\begin{equation}
  \frac{d\sigma}{d\pT} = \frac{N_{\text{photon}}}{\Delta\pT \times \Lint \times \varepsilon_{\text{total}}}
\end{equation}
with
\begin{equation}
  \varepsilon_{\text{total}} = \varepsilon_{\text{reco}} \times
  \varepsilon_{\text{iso}} \times \varepsilon_{\text{id}} \times \veff
\end{equation}

Two equivalent conventions are possible for $\Lint$ and $\veff$:

\textbf{Old convention (used in production until now):}
\begin{equation}
  \Lint = \Lint^{60\text{cm}}, \qquad
  \veff^{\text{old}} = \frac{N(\text{truth }\gamma,\ |z_t|<60,\ \text{MBD+}|z_r|<60)}{N(\text{truth }\gamma,\ |z_t|<60)}
\end{equation}
with $\Lint^{60\text{cm}}$ from \texttt{60cmLumi\_fromJoey.list},
computed by \texttt{get\_luminosity\_182630.C(..., zsel=1, ...)}
with the scaler tally restricted to $|\zvtx|<60$~cm.

\textbf{New convention:}
\begin{equation}
  \Lint = \Lint^{\text{noVtx}}, \qquad
  \veff^{\text{new}} = \frac{N(\text{truth }\gamma,\ |z_t|<60,\ \text{MBD+}|z_r|<60)}{N(\text{truth }\gamma,\ \text{no vertex cut})}
\end{equation}
with $\Lint^{\text{noVtx}}$ from \texttt{allzLumi\_fromJoey.list}
(\texttt{zsel=5}, no vertex cut at the scaler level).

The numerators are identical. Only the denominator of the efficiency
factor changes (and the paired luminosity). The new convention is
preferred because it cleanly separates the physics fiducial
cross-section (all truth photons in the kinematic region) from the
detector acceptance (MBD + vertex).

The two conventions must give the same final cross-section to
$\leq 1\%$ in the single-interaction limit; any larger deviation
indicates a bug in the MC vertex reweighting (which is the one
thing that distinguishes the two denominators).

%% ============================================================
\section{The Bug in the Old Production}
\label{sec:bug}

Using the production \texttt{MC\_efficiency\_nom.root} (June 2025),
the two conventions give wildly different cross-sections:

\begin{table}[H]
\centering
\caption{Pre-fix comparison: same numerator, different denominators.
  If the MC vertex distribution matched data, the ratio would be
  $\approx 1$. Instead the new convention gives a product $\sim 2.8$
  times smaller.}
\label{tab:prefix}
\small
\begin{tabular}{lrrrrr}
\toprule
\textbf{Period} & \textbf{$\veff^{\text{old}}$} & \textbf{$\veff^{\text{new}}$}
  & \textbf{$L_{60\text{cm}}\!\cdot\!\veff^{\text{old}}$}
  & \textbf{$L_{\text{noVtx}}\!\cdot\!\veff^{\text{new}}$}
  & \textbf{Ratio} \\
\midrule
1.5 mrad & 0.676 & 0.240 & 11.41 & 4.12  & 0.361 \\
0 mrad   & 0.676 & 0.240 & 22.10 & 11.34 & 0.513 \\
\bottomrule
\end{tabular}
\end{table}

The diagnosis: the old \texttt{RecoEffCalculator\_TTreeReader.C}
used a static fallback reweight file
(\texttt{results/vertex\_reweight\_bdt\_none.root}) when the
on-the-fly path failed. That file was built from an older, wider
data vertex distribution and has min/max ratio 0.49 / 1.78 ---
nearly flat. At large $|\zvtx|$, MC events were kept at $\sim 50\%$
weight instead of being suppressed. The fill at
\texttt{RecoEffCalculator\_TTreeReader.C:1691} (no-vertex-cut
$h_{\text{truth}\ \pT}$) picked up these bloated events, so the
denominator of $\veff^{\text{new}}$ was $\sim 2.8$ times too large.

A further subtlety: the lookup at the event level had a fallback
\texttt{vertex\_weight = 1.0} for non-finite or $\leq 0$ bin contents.
This kept raw MC events outside the data acceptance at full weight,
amplifying the bug.

%% ============================================================
\section{The Fix}
\label{sec:fix}

Three changes were applied in commits 8ac4d1b and 8ecae87 to
\texttt{RecoEffCalculator\_TTreeReader.C} and
\texttt{ShowerShapeCheck.C}:

\begin{enumerate}
  \item \textbf{Removed the static fallback} entirely. If vtxscan
    files are missing, the code now prints \texttt{FATAL} and
    returns. The static file was renamed to \texttt{*.deprecated}
    on disk.
  \item \textbf{Hardened the on-the-fly reweight}:
  \begin{itemize}
    \item Rebin \texttt{h\_vertexz} (default factor 5 $\rightarrow$
      5~cm bins, configurable via
      \texttt{vertex\_reweight\_rebin}) to improve tail statistics.
    \item Smooth the data/sim ratio with \texttt{TH1::Smooth()}
      (default 1 pass, configurable via
      \texttt{vertex\_reweight\_smooth}).
    \item Linear interpolation between bin centers with
      \texttt{TH1::Interpolate()} for smooth event-level weights.
    \item \texttt{vertex\_weight = 0.0} fallback for out-of-range
      or non-finite values. Previously \texttt{1.0}, which kept raw
      MC events at large $|\zvtx|$.
  \end{itemize}
  \item \textbf{Added the new convention as a config-switchable
    path} in \texttt{CalculatePhotonYield.C}
    (\texttt{use\_lumi\_novtx: 1}). The old convention remains the
    default. Both paths read from the same MC output; only the
    denominator histogram and luminosity value change.
\end{enumerate}

The physics motivation for the zero fallback: if data has
essentially no events at some $|\zvtx|$, MC events at that position
do not contribute to the observed cross-section and should be
dropped. Keeping them at weight 1.0 is equivalent to claiming those
MC events represent data, which is wrong.

%% ============================================================
\section{Post-Fix Validation}
\label{sec:validation}

A standalone simulation
(\texttt{efficiencytool/lumi\_convention\_validation.py}) applies the
fixed reweight logic to the existing first-pass vtxscan files
(\texttt{data\_histo\_bdt\_nom\_vtxscan.root} and
\texttt{MC\_efficiency\_photon10\_bdt\_nom\_vtxscan.root}) and
propagates the result through the cross-section formula. This
simulates what the fixed production pipeline would produce without
actually re-running MergeSim.

\subsection{Vertex reweight shape}

Figure~\ref{fig:reweight_shape} compares the old and new reweight
functions. The two are very similar in the bulk of the distribution
($|\zvtx| < 50$~cm) --- the bug only affects the tails, where the
old code's \texttt{fallback = 1.0} kicks in for bins with
poorly-defined ratio.

\begin{figure}[H]
\centering
\includegraphics[width=0.80\textwidth]{vertex_reweight_shape.pdf}
\caption{Vertex reweight function (data/sim, normalized) computed
  from the 1.5~mrad nominal vtxscan files. Old (blue) uses 1~cm
  bins with \texttt{fallback = 1.0}; new (red) uses 5~cm rebin,
  one smoothing pass, and \texttt{fallback = 0.0}. The two agree
  in the bulk; the tail behavior is the diagnostic region where
  the old code's bug lived.}
\label{fig:reweight_shape}
\end{figure}

\subsection{Reweighted MC vertex distribution vs.\ data}

Figure~\ref{fig:vertex_dist} overlays the raw MC, the old and new
reweighted MC, and the data vertex distributions. Both reweighted
distributions match the data shape very well.

\begin{figure}[H]
\centering
\includegraphics[width=\textwidth]{vertex_dist_comparison.pdf}
\caption{Reconstructed vertex $z$ distribution from data (black),
  raw MC (green dashed), old-reweighted MC (blue), and
  new-reweighted MC (red). Left: linear scale. Right: log scale.
  The dotted vertical lines mark $|\zvtx|=60$~cm. Both old and new
  reweights reproduce the data shape at the histogram-integral
  level; the bug was in how event-level weights were used
  downstream.}
\label{fig:vertex_dist}
\end{figure}

\subsection{Integrated vertex acceptance}

After reweighting, the integrated MC fraction within $|\zvtx|<60$~cm
should match the data fraction of 0.9865 (1.5~mrad target).

\begin{table}[H]
\centering
\caption{Integrated $|\zvtx|<60$ fraction from the 1.5~mrad vtxscan
  files.}
\label{tab:integral}
\begin{tabular}{lr}
\toprule
\textbf{Source} & \textbf{Fraction $|\zvtx|<60$} \\
\midrule
Data (target)          & 0.9865 \\
Raw MC                 & 0.7296 \\
Old-reweighted MC      & 0.9865 \\
New-reweighted MC      & 0.9862 \\
\bottomrule
\end{tabular}
\end{table}

Both old and new reweights reproduce the data fraction to
$\leq 0.03\%$. The pre-fix production's 0.355 ratio came from the
\emph{static fallback file} being used (not the on-the-fly path),
which has min ratio 0.49 so MC events at large $|\zvtx|$ were kept
at $\sim 50\%$ weight.

\subsection{Per-\texorpdfstring{$\pT$}{pT} consistency check}

The key cross-check: after the fix, does
$\Lint^{\text{noVtx}} \times \veff^{\text{new}}$ agree with
$\Lint^{60\text{cm}} \times \veff^{\text{old}}$ bin-by-bin?

\begin{table}[H]
\centering
\caption{Post-fix cross-check at 1.5~mrad. $\veff^{\text{new}}$ is
  computed assuming the fixed MC pipeline would produce
  $h_{\text{truth}\ \pT}^{\text{new}}$ with integral equal to
  $h_{\text{truth}\ \pT\ \text{vertexcut}} / 0.9865$. The per-bin
  ratio is flat at $\sim 1.003$ across all \pT\ bins.}
\label{tab:perbin}
\small
\begin{tabular}{lrrrrr}
\toprule
\textbf{\pT\ [GeV]} & \textbf{$\veff^{\text{old}}$}
  & \textbf{$\veff^{\text{new}}$}
  & \textbf{$L_{60} \cdot \veff^{\text{old}}$}
  & \textbf{$L_{\text{nov}} \cdot \veff^{\text{new}}$}
  & \textbf{Ratio} \\
\midrule
8--10    & 0.6771 & 0.6680 & 11.43 & 11.47 & 1.0032 \\
10--12   & 0.6670 & 0.6580 & 11.26 & 11.29 & 1.0032 \\
12--14   & 0.6568 & 0.6479 & 11.09 & 11.12 & 1.0032 \\
14--16   & 0.6452 & 0.6365 & 10.89 & 10.92 & 1.0032 \\
16--18   & 0.6347 & 0.6261 & 10.71 & 10.75 & 1.0032 \\
18--20   & 0.6219 & 0.6135 & 10.50 & 10.53 & 1.0032 \\
20--22   & 0.6095 & 0.6012 & 10.29 & 10.32 & 1.0032 \\
22--24   & 0.6027 & 0.5946 & 10.17 & 10.21 & 1.0032 \\
24--26   & 0.5789 & 0.5711 &  9.77 &  9.80 & 1.0032 \\
26--28   & 0.5649 & 0.5573 &  9.53 &  9.56 & 1.0032 \\
32--36   & 0.4999 & 0.4931 &  8.44 &  8.46 & 1.0032 \\
\midrule
Integrated & 0.6758 & 0.6667 & 11.41 & 11.44 & 1.0032 \\
\bottomrule
\end{tabular}
\end{table}

Figure~\ref{fig:ratio} shows the same information graphically:
both conventions give the same $\Lint \times \veff$ product within
$+0.32\%$, flat across all \pT\ bins. A shift this small is well
below any relevant PPG12 systematic (purity $\sim 5$--$10\%$, energy
scale $\sim 3$--$5\%$, luminosity $\sim 7\%$) and is consistent
with zero within the statistical precision of the vtxscan files
used to simulate the fix.

\begin{figure}[H]
\centering
\includegraphics[width=0.75\textwidth]{L_times_eff_ratio.pdf}
\caption{Top: $\Lint \times \veff$ per \pT\ bin for the old
  convention ($L_{60\text{cm}} \times \veff^{\text{old}}$, blue)
  and the new convention
  ($L_{\text{noVtx}} \times \veff^{\text{new}}$, red). Bottom:
  ratio of new to old. Flat at $+0.32\%$ across all \pT\ bins.}
\label{fig:ratio}
\end{figure}

\subsection{Rebin stability}

Figure~\ref{fig:rebin_stability} shows the fixed reweight computed
with different rebin factors. The shape is stable from 1~cm (no
rebin) to 20~cm bins in the bulk of the distribution. The default
choice of 5~cm bins gives smooth coverage in the tails without
over-coarsening the peak.

\begin{figure}[H]
\centering
\includegraphics[width=0.80\textwidth]{rebin_stability.pdf}
\caption{Vertex reweight function computed with rebin factors
  1, 2, 5, 10, 20 (one smoothing pass each). The shape is stable
  across a factor of 20 in bin width.}
\label{fig:rebin_stability}
\end{figure}

%% ============================================================
\section{Implementation}
\label{sec:implementation}

To enable the new convention in any analysis config YAML, set:

\begin{verbatim}
analysis:
  use_lumi_novtx: 1
  lumi: 17.1642   # 1.5 mrad, from allzLumi_fromJoey.list
  # use 47.2284 for 0 mrad, 64.3718 for combined
\end{verbatim}

The default is \texttt{use\_lumi\_novtx: 0} (old convention). The
original \texttt{vertex\_eff} calculation in
\texttt{CalculatePhotonYield.C} is kept as the default path; the
new path is an \texttt{if} branch inside the efficiency loop
(lines 1063--1071).

New config knobs for the vertex reweight (both have sensible
defaults, so existing configs work unchanged):

\begin{verbatim}
analysis:
  vertex_reweight_rebin: 5       # rebin factor for h_vertexz (1cm → 5cm)
  vertex_reweight_smooth: 1      # TH1::Smooth() passes
\end{verbatim}

The old \texttt{vertex\_reweight\_file} config knob is no longer
read (the static fallback path is removed).

%% ============================================================
\section{Caveats and Remaining Work}
\label{sec:caveats}

\begin{enumerate}
  \item \textbf{The validation uses a simulated post-fix result},
    not an actual re-run. The assumption is that
    $h_{\text{truth}\ \pT}^{\text{new}}$ integral equals
    $h_{\text{truth}\ \pT\ \text{vertexcut}}$ integral divided by
    $0.9865$, based on the shape-level vertex reweight validation.
    This is a good approximation because the vertex reweight within
    $|\zvtx|<60$~cm is unchanged by the fix. A true re-run
    (\texttt{oneforall.sh config\_bdt\_nom.yaml}) is recommended to
    confirm the $+0.32\%$ shift holds in the actual production MC.

  \item \textbf{Per-segment vertex variation}: the current reweight
    uses a single global vertex distribution per analysis config.
    Within a fill, the luminous region moves as the beam decays,
    but this effect is typically $\leq 1\%$ at RHIC luminosities.

  \item \textbf{Pileup}: this note does not modify the pileup
    correction. The double-interaction MC blending in
    \texttt{run\_showershape\_double.sh} is unaffected and carries
    its own correction.

  \item \textbf{0~mrad consistency}: the validation was done with
    the 1.5~mrad vtxscan files. The 0~mrad lumi has a larger
    $L_{\text{noVtx}} / L_{60\text{cm}} = 1.44$ ratio, reflecting a
    much wider vertex distribution. The fix should work there too,
    but the actual post-fix result should be verified when the
    0~mrad config is re-run.

  \item \textbf{Re-run is cheap}: the fix does not require
    regenerating vtxscan files. Only the second pass of
    \texttt{oneforall.sh} needs to re-read the slimtree. Once the
    re-run completes, the new \texttt{MC\_efficiency\_nom.root}
    should have
    $h_{\text{truth vertexcut}}.\text{Integral()} /
    h_{\text{truth}}.\text{Integral()} \approx 0.9865$, and the
    \texttt{use\_lumi\_novtx: 1} path should give a cross-section
    within $\sim 0.3\%$ of the \texttt{use\_lumi\_novtx: 0} path.
\end{enumerate}

%% ============================================================
\section{Files}
\label{sec:files}

\begin{tabular}{ll}
\toprule
\textbf{File} & \textbf{Change} \\
\midrule
\texttt{efficiencytool/RecoEffCalculator\_TTreeReader.C} & Fix vertex reweight, remove fallback \\
\texttt{efficiencytool/ShowerShapeCheck.C}               & Same fix \\
\texttt{efficiencytool/CalculatePhotonYield.C}           & Add \texttt{use\_lumi\_novtx} path \\
\texttt{efficiencytool/lumi\_convention\_validation.py}  & Validation script \\
\texttt{plotting/plot\_lumi\_convention.py}              & Plotting script \\
\texttt{lumi/allzLumi\_fromJoey.list}                    & New: no-vertex-cut lumi \\
\texttt{results/vertex\_reweight\_bdt\_none.root.deprecated} & Deprecated fallback \\
\bottomrule
\end{tabular}

\end{document}
